\documentclass[11pt]{article}
\usepackage[T1]{fontenc}
\usepackage{amsmath,amssymb,amsthm}
\usepackage[margin=1in]{geometry}
\usepackage{xurl}
\usepackage[hidelinks]{hyperref}
\newtheorem{proposition}{Proposition}
\newtheorem{lemma}[proposition]{Lemma}
\newcommand{\Tr}{\operatorname{Tr}}
\newcommand{\Ann}{\operatorname{Ann}}
\newcommand{\Spec}{\operatorname{Spec}}
\newcommand{\kerop}{\operatorname{ker}}
\title{The residue-field case: reducing the proof to what the local theorem uses}
\author{}
\date{31 October 2026}
\begin{document}
\maketitle

\noindent\textbf{Status.}
The diagonal, general flatness transfer and field regularity in
Sections~2--5 are now checked in Unit, including the finite span, an
actual finite basis and a perfect trace pairing. The finite-field
primitive-element theorem, the small-coefficient-field branch, selected
relative bases and selected simple-root equations also check. The
large-field diagonal search and both field endpoints now check:
finite presentation and formal unramifiedness give a finite product of
separable monic quotients, without a prime parameter, as well as a
separable monogenic neighbourhood of a specified decidable prime.
The product proof simplifies the large-field argument by using a
divided difference and trace duality directly. The remaining local work
is the connection supplying all local-base inputs.
The starting point is Lombardi's
\emph{Un th\'eor\`eme d\'ecisif} \cite{decisif} and Lombardi--Quitt\'e
\cite{LQ}.  The suggested reorganisation uses the flatness and
idempotent-kernel arguments already checked in
\path{sag/etale-flat-unramified}.

Its transfer theorem for a standard
\emph{\'etale} source has now been generalised to a finitely presented,
formally unramified source in \path{sag/etale-field-transfer}.
The diagonal construction is in \path{sag/etale-field-diagonal};
the field application is in \path{sag/etale-field-regularity}.
The combined theorem is \texttt{UnramifiedFieldStructure} in
\path{sag/etale-field}; no finite basis or intermediate theorem is
assumed in its input. The selected-prime conclusion is
\texttt{UnramifiedFieldSelected} in \path{sag/etale-field-complete};
it constructs the basis, trace pairing, diagonal and primitive factor
from those same hypotheses and the specified prime.
The full product conclusion is \texttt{UnramifiedFieldProduct} in
\path{sag/etale-field-product}; it constructs both directions of the
algebra isomorphism and uses no specified prime or factorisation oracle.

\section{Start from the two actual outputs}

Let $k$ be a nonzero discrete field and $C$ a finitely presented
$k$-algebra, unramified in the sense of uniqueness of lifting across
square-zero ideals.

First, the ZMT application needs that $C$ is finite over $k$.
A finite basis is a useful stronger output, provided it is constructed,
not inferred from finite generation alone.

Second, the local structure proof has a \emph{specified decidable prime}
$\mathfrak p$ of $C$.  It needs only a monogenic separable factor which
contains that prime.  A sufficient field-level output is
\[
 e^2=e,\qquad e\notin\mathfrak p,\qquad
 C_e=k[x],\qquad P(x)=0,\qquad P'(x)\in C_e^\times,
\]
with $P\in k[T]$ monic.  An isomorphism with $k[T]/(P)$ is more than
the downstream proof requires; a separable annihilating polynomial
and generation suffice.  Nor does it need a decomposition into fields,
irreducibility tests, or all factors at once.

There is a further finite-interface simplification.  The current
\texttt{Introduction.Generates} asks for polynomial representations
of every $b\in B$, with untruncated witnesses.  The only actual
consumption, after the change $x\mapsto x-1$, is at the finitely many
module generators \texttt{bs i}: see \texttt{Introduction.seg}, the
declarations \texttt{Ci} and \texttt{e-gen}.  The field bridge should
therefore supply merely the representations of these finitely many
elements.  Finite choice then suffices. This refactoring now checks in
\path{sag/etale-field-local-interface}: private variants of the introduction
and local proof consume just these finitely many representations.
\texttt{SelectedFlatLocal} takes one factor outside the specified prime,
without requiring a global product decomposition. The earlier files
are unchanged.

\section{A common certificate: the diagonal idempotent}

\begin{lemma}\label{diagonal}
Let $R\to S$ be a finitely presented, formally unramified algebra.
Put $D=S\otimes_R S$ and $J=\kerop(\mu:D\to S)$.
There exists $\epsilon\in D$ such that
\[
 \mu(\epsilon)=1,\qquad
 (s\otimes1)\epsilon=(1\otimes s)\epsilon\quad(s\in S).
\]
Moreover $\epsilon^2=\epsilon$.
\end{lemma}
\begin{proof}
If $x_1,\ldots,x_n$ generate $S$, the differences
$x_i\otimes1-1\otimes x_i$ generate $J$.
Consider the two maps $S\rightrightarrows D/J^2$.
They become equal in $D/J\simeq S$; the kernel $J/J^2$ of this last
quotient is square-zero.  Uniqueness of lifting makes the two maps
equal already in $D/J^2$.  Thus each difference belongs to $J^2$ and
$J=J^2$.

Apply the determinant lemma for a finitely generated idempotent ideal.
It supplies $\epsilon$ with $\epsilon J=0$ and $1-\epsilon\in J$.
These are precisely the two required identities, and
$\epsilon(1-\epsilon)=0$ gives idempotence.
\end{proof}

This proof starts from the existing \texttt{FormallyEtale.Unique}
interface.  It needs no construction of the full module of
K\"ahler differentials or an equivalence with the Jacobian criterion.
The quotient and determinant infrastructure already exists.  The
underlying diagonal argument is also recorded in \cite{diagonal}.

\section{Extend the checked flatness argument}

\begin{proposition}\label{transfer}
Suppose $R\to S$ has the certificate of Lemma~\ref{diagonal}.
For any $R$-algebra map $S\to T$, if $T$ is flat over $R$, it is flat
over $S$.
\end{proposition}
\begin{proof}
The algebra $D_T=S\otimes_R T$ is flat over $S$ by flat base change.
Map the second factor of $\epsilon$ to $T$, obtaining $\epsilon_T$.
The multiplication map $\mu_T:D_T\to T$ has an $S$-linear section
\[
 t\longmapsto\epsilon_T(1\otimes t).
\]
Linearity follows from the equalisation identity, and the composite
with $\mu_T$ is the identity since $\mu_T(\epsilon_T)=1$.
A module retract of a flat module is flat.
\end{proof}

The proof is the same averaging mechanism used in the checked standard
\'etale transfer theorem.  There, synthetic division constructs the
averaging element.  Here, formal unramifiedness and the determinant
lemma construct it.

Both extensions now check.  \texttt{DiagonalCertificate} constructs
the diagonal certificate from \texttt{Unique}.
\texttt{FlatPresentedCocone} iterates polynomial base change and
simultaneous principal quotients, then transports the source along
its presentation isomorphism.  It supplies a flat algebra with the two
structure maps and the mapping-out laws needed for averaging; it does
not require a finite presentation of $T$.
\texttt{FlatDiagonalTransfer} combines these constructions.

\section{Finiteness without the derivation argument}

Throughout this note, \emph{zero-dimensional} means the algebraic
property
\[
 \forall x\;\exists N\in\mathbb N,\ b\qquad x^N(1-xb)=0.
\]
Existence is propositional truncation in Unit.  No definition through
prime ideals is used, and $N=0$ is allowed.

\begin{lemma}\label{regular}
For $C$ as above and every $a\in C$, there is merely a $b\in C$ with
$a=a^2b$.  Consequently $C$ is reduced and zero-dimensional.
\end{lemma}
\begin{proof}
Every $k$-algebra is $k$-flat by the equational criterion.  This does
not require a basis: for a finite relation $\sum\alpha_i t_i=0$,
either all $\alpha_i$ vanish or one is invertible.  In the second case
the columns $e_i-(\alpha_i/\alpha_p)e_p$ explain the relation.

Apply Proposition~\ref{transfer} to $C\to C/(a)$.  This quotient is
$C$-flat.  The checked flat-kernel lemma gives $a=az$ with $z\in(a)$;
writing $z=ab$ yields $a=a^2b$.
This is von Neumann regularity, which implies both assertions.  For
example, if $a^m=0$, iterating $a=a^2b$ gives $a=a^m b^{m-1}=0$.
\end{proof}

This removes the rational-function derivation extension from
Lombardi's Theorem~4.5.  It also removes primitive elements, perfection
and discriminants from the \emph{finiteness} stage.  It uses neither
ZMT nor the theorem flat plus unramified implies \'etale.  Its dependencies
are the explicit diagonal, base-change, and retraction arguments above.

The finite-span step can be simplified further than the initial version
of this note.  Induction on \emph{algebra generators}, rather than on a
chosen finite presentation, avoids contraction ideals entirely.

\begin{proposition}\label{finite}
A finitely generated zero-dimensional algebra over a discrete field
is finite.  If it is finitely presented, it also has a finite basis.
\end{proposition}
\begin{proof}[Construction]
Induct on the number $n$ of algebra generators $x_1,\ldots,x_n$ of $C$.
For $n=0$, the single element $1$ spans $C$ over $k$, including when
$C$ is the zero algebra.  No decision between $C=0$ and $C=k$ is needed.

For $n>0$, zero-dimensionality supplies
$x_1^N(1-x_1b)=0$.  Write $b=P(x_1,\ldots,x_n)$, using algebra
generation.  The polynomial
\[
 f=X_1^N(1-X_1P)
\]
vanishes in $C$ and is nonzero: its coefficient at
$X_1^NX_2^0\cdots X_n^0$ is $1$.  This also works when $N=0$.
A nonzero constant relation gives the zero algebra, which is finite.
Otherwise a triangular change of variables of Nagata type makes $f$
monic in the last variable $T$, after dividing by its nonzero leading
scalar.  With every exponent less than $M$, the substitutions
$X_i=Y_i+T^{M^i}$ for $i<n$, keeping $X_n=T$ fixed, give distinct top
$T$-degrees to the original monomials.  This works over small fields.

Let $A'\subseteq C$ be the image subalgebra generated by the images of
$Y_1,\ldots,Y_{n-1}$, and let $t$ be the image of $T$.  The inverse
coordinate substitution shows $C=A'[t]$.  The monic equation of degree
$d$ makes $1,t,\ldots,t^{d-1}$ span $C$ over $A'$, so $C$ is integral
over $A'$ by the determinant trick.

Zero-dimensionality descends along this integral injection.  Explicitly,
for $x\in A'$, suppose
\[
 x^N=b x^{N+1},\qquad
 b^h+c_{h-1}b^{h-1}+\cdots+c_0=0,
\]
where the $c_i$ belong to $A'$.  The elementary \emph{power absorption}
lemma, proved by induction, is
\[
 x^N=b^j x^{N+j}\qquad(j\geq0).
\]
Consequently $b^j x^{N+h}=x^{N+h-j}$ for $0\leq j\leq h$.
Multiplying the monic equation by $x^{N+h}$ therefore gives
\[
 x^N(1+c_{h-1}x+\cdots+c_0x^h)=0.
\]
Injectivity makes this an equation in $A'$.  That subalgebra is generated
by $n-1$ elements, so induction makes it finite over $k$.  Combining
the two finite spanning families proves that $C$ is finite over $k$.
Neither a finite presentation of $A'$ nor generators of a contraction
ideal are needed.

For the additional basis conclusion, use the finite presentation
\[
 C=k[X_1,\ldots,X_n]/(f_1,\ldots,f_m).
\]
To obtain a basis, do not merely assume that a finitely generated
$k$-module has one.  Use the finite spanning family and the determinant
trick to obtain monic annihilating polynomials $p_i(X_i)$ for the
original generators.  Then
\[
 Q_0=k[X_1,\ldots,X_n]/(p_1(X_1),\ldots,p_n(X_n))
\]
has its explicit bounded-monomial basis.  The original equations
$f_i$ generate an ideal of $Q_0$ with finite $k$-spanning family
consisting of their products with those basis monomials.  Gaussian
elimination on this finite matrix constructs a basis of the quotient
$C$.  This last step uses the original finite presentation essentially.
\end{proof}

The earlier contraction route remains valid: coherence of a finite free
polynomial module gives generators of a contraction by projecting the
last coordinates of the syzygies of $[V,-e_0]$.  It is unnecessary in
the simplified finite-span proof above. The primitive-element stage is
now checked as described below. The Nagata substitution, generated
image subalgebra, integral descent and finite-span induction now all
check in \path{sag/etale-field-finite} and
\path{sag/etale-field-generated}. Their combined theorem is
\texttt{ZeroDimensionalFinite}.

The application
\texttt{UnramifiedField\allowbreak FiniteGeneral} takes just finite presentation and
formal unramifiedness over a discrete nonzero field.

The separate basis assembly also checks as
\texttt{FinitePresented\allowbreak FieldBasis}. It constructs the iterated monic
cover, proves that the original equations generate its kernel, and
forms the finite matrix consumed by quotient Gaussian elimination.
Neither that matrix nor the cover's universal property is assumed.
The file named \texttt{WeakNullstellensatz} in the affine
tree is an SQC statement; it is not this finiteness theorem.

\section{A small trace identity after finiteness}

This implication now checks as \texttt{FiniteUnramifiedTrace} in
\path{sag/etale-field-trace}, with an explicit finite basis as input.
The scalar-extension algebra and its coordinates in the algebra tensor
product are constructed; no tensor-basis hypothesis is assumed.

With a basis $b_1,\ldots,b_d$ of $C$, write the diagonal element as
$\epsilon=\sum_i b_i\otimes c_i$.
Let $\lambda_i$ be the coordinate functions.  Comparing coefficients
in the equalisation identity gives, for every $z\in C$,
\[
 zc_i=\sum_j\lambda_i(zb_j)c_j.
\]
Put $z=ab_i$ in the $i$th identity and sum over $i$.  The decisive
coordinate calculation is
\[
 \sum_i\lambda_i(ab_i b_j)
   =\sum_i\lambda_i((ab_j)b_i)=\Tr_{C/k}(ab_j),
\]
since these are the diagonal entries of multiplication by $ab_j$.
Since $\sum_i b_i c_i=1$, exchanging the two finite sums gives
\[
 a=\sum_j\Tr_{C/k}(ab_j)c_j.
\]
Hence the trace-pairing matrix has an inverse; its determinant is a
unit.  This is the $\mu(\epsilon)=1$ specialisation of the trace
identity in Quitt\'e--Lombardi \cite[Theorem~3.1]{trace}.
Only this specialisation is needed, not the full norm--different
divisibility theorem.  The diagonal certificate has also avoided the
detour through a perfect closure in \cite[Theorem~4.3]{decisif}.

\section{Selected linear algebra and simple roots}

The relative-basis step can be confined to one specified decidable prime
of a regular ring $A$. Start with a finite presentation
$A^m\to A^n\to M\to0$. The quotient $A/\mathfrak p$ is a discrete
field: a regular witness $a=a^2b$ gives the inverse of $\bar a$ when
$a\notin\mathfrak p$. Run ordinary Gaussian elimination there, and lift
only its finite projection, section and relation-coefficient matrices.
Their equation errors form a finite family in $\mathfrak p$.

For an error $a\in\mathfrak p$, a regular witness supplies the idempotent
$1-ab\notin\mathfrak p$ which kills $a$. The product of these finitely
many idempotents kills every matrix error. Thus on one selected factor
$eA$, the lifted matrices give exact quotient coordinates and an actual
basis of $eM$. This construction is checked in
\path{sag/etale-field-selected-gaussian}; its algebra-basis adapter is
\path{sag/etale-field-selected-basis}. No full decomposition of $A$
is constructed.

There is also a smaller replacement for the discriminant calculation
which makes a monogenic polynomial separable. For a finite algebra $A$
with nondegenerate trace and a $k$-derivation $D:A\to A$, the commutator
identity
\[
 [bD,m_a]=m_{bD(a)}
\]
gives $\Tr(bD(a))=0$ for every $b$, hence $D(a)=0$.
For $Q=k[T]/(f)$, an annihilator $h$ of $f'(T)$ defines the root
$T+h\varepsilon$ in the dual numbers and therefore an actual
derivation with $D(T)=h$.

Only trace nondegeneracy on a selected factor $eQ$ is needed. Construct
the relative derivation $Q\to eQ$ and restrict it to $eQ$.
The equation $D(e)=D(e)+D(e)$ gives $D(e)=0$, without dividing by $2$,
and $D(eT)=D(T)=h$. Nondegeneracy makes $h=0$.
Finally regularity $f'=(f')^2b$ and the vanishing of the annihilator
give $f'b=1$ on this factor.

These steps now check. For an element $x\in C$, the constructed minimal
quotient $Q=k[T]/(f)$ embeds in $C$. Selected relative bases, rectangular
trace cyclicity, and trace transitivity pass nondegeneracy first to the
selected factor of $C$ and then to $eQ$. The result in
\path{sag/etale-field-selected-separable} is a monic polynomial which
annihilates $x$ globally and has invertible derivative on a factor
outside the specified prime. For a finite family of $x$'s, multiply
their independently constructed idempotents; every derivative inverse
equation remains valid. This does not assert that the original
polynomials are separable on their discarded factors.

\section{The selected primitive-element proof}

LQ VI--1.4, VI--1.7 and VI--1.9 give separable monogenic factors from
the finite basis and invertible trace pairing. The preceding selected
trace argument supplies the simple-root part. A further simplification
of the large-field choice uses the diagonal directly, as
described below; no irreducible factorisation over an arbitrary discrete
field is being assumed.

For the local theorem, carry the specified prime through each
idempotent refinement. Decidability and primality select the side
containing the prime; products of these idempotents retain it.
The checked construction does not retain the discarded factors or
assemble a global product decomposition.

The small-field branch remains necessary.  For example
$\mathbf F_2^3$ is unramified and finite but is not monogenic:
every element satisfies $X^2-X$, so the subalgebra it generates has
dimension at most two.  At a specified prime, however, one coordinate
factor is $\mathbf F_2$ and is monogenic.

LQ's large-field argument chooses a good linear combination by a
nonzero discriminant polynomial.  Its alternative branch works over
the finite field generated by the relevant coefficients.  A useful
implementation of this dichotomy is bounded closure: grow a finite
set containing $0,1$ and those coefficients, closing under addition,
negation, multiplication and total inverse (with $0^{-1}=0$), until it
is a finite subfield or has more elements than the required bound.
Including inverses avoids a separate finite-injection argument for the
subfield. This construction now checks as \texttt{FiniteCoefficientField};
it does not decide whether the whole field is finite.

The finite-field primitive-element argument is now checked as well.
Choose an element of maximal multiplicative order by finite search.
An explicit gcd-stripping construction combines two orders into an
order divisible by both. Maximality then makes the selected order a
period of every unit. Its distinct powers, together with any alleged
missing nonzero element, would give too many roots of $T^m-1$.
The polynomial root bound excludes that case. The common monic
annihilator $T^q-T$ has derivative $-1$; both assertions are proved
from the finite enumeration. This uses no field classification or
irreducible factorisation.

The small-field assembly constructs the finite subalgebra generated
over the actual finite coefficient field, selects its field factor,
and applies this primitive-element theorem. Its inclusion in the
ambient algebra sends $1$ to the selected idempotent $e$.
The checked \texttt{CornerHorner} identities retain that factor in
both values and derivatives. The resulting \texttt{SmallFieldBranch}
produces exactly the selected certificate required by the local proof,
from the finite coefficient field and the finitely many monic equations.
These equations are themselves constructed by the determinant trick
from the finite basis. Their lower coefficients are collected in one
finite family before applying bounded closure. In the small branch,
each occurrence has an explicit lift to the resulting finite subfield,
so the actual monic lists and their root equations are available there.

\subsection*{A diagonal route for the large-field choice}

Here is the checked replacement for the universal splitting-algebra
and discriminant calculation. Write $D=C\otimes_k C$ and
$\Delta x=x\otimes1-1\otimes x$. For a basis $b_1,\ldots,b_n$,
the diagonal certificate implies that $\epsilon$ and the
$\Delta b_i$ generate the unit ideal of $D$. Consider
\[
 g(T)=\epsilon+\sum_{i=1}^n\Delta b_i T^i\in D[T].
\]
Its coefficients generate the unit ideal. Polynomial McCoy therefore
makes multiplication by $g$ injective. On the constructed finite
$k[T]$-basis of $D[T]$, matrix McCoy makes its determinant a nonzero
polynomial. Its degree is at most $n\dim_k D=n^3$.
Sufficiently many distinct scalars and the root bound supply an $a$
where the determinant is nonzero, making $g(a)$ invertible.
For $z=\sum_i a^i b_i$, this gives
\[
 (\Delta z)w=1-\epsilon.
\]
The determinant, polynomial-basis and specialisation assembly for this
search is checked in \path{sag/etale-field-diagonal-separator}.
It requires $n^3+1$ distinct scalars and constructs both $z$ and $w$.

The implication from this certificate to the needed coordinate
equations is checked as well. If $p:C\to T$ is an algebra map and a
$k$-linear map $\lambda:C\to T$ satisfies
$\lambda(zc)=p(z)\lambda(c)$, contract the left tensor factor by
$\lambda$. This kills every multiple of $\Delta z$. Since
\[
 \Delta c=(\Delta z)(w\Delta c),
\]
it gives $\lambda(c)=\lambda(1)p(c)$.
Apply this to the coordinates of a selected relative basis over
$k[z]$. Expanding $1$ selects one coordinate $\lambda_j(1)$ outside
the prime. Regularity makes it invertible on a further selected factor,
where the displayed equation expresses every element as an element
of $k[z]$. The coordinate and monic-representative constructions now
check; they avoid a separate relative-rank-one argument.

The final assembly in \path{sag/etale-field-selected-complete} applies
bounded coefficient closure with this scalar bound. Its two branches
return the same selected certificate. In the large branch the
generation and simple-root idempotents are intersected; in the small
branch the finite subfield and its lifted monic equations feed the
finite-field construction. All finite choices remain under truncation.
Neither a large-field hypothesis nor a decision that the whole field
is finite is required.

\section{The full product theorem and a shorter large-field argument}

The checked endpoint is now
\begin{center}
\texttt{UnramifiedFieldProduct.Over[k,setk,field]}\\
\texttt{.At[C,fp,unique].theorem}.
\end{center}
It merely constructs
\[
 C\simeq\prod_i k[T]/(f_i),
 \qquad f_i\text{ monic},\quad f_i'\text{ invertible modulo }f_i.
\]
Both directions are algebra maps with proved inverse equations.
The finite product is represented by an ordinary finite binary tree.
Zero factors are allowed, represented by the monic polynomial $1$.
No irreducibility is asserted: these are separable monogenic algebras,
and decomposition into fields would require additional factorisation
information about $k$.

\subsection{A tensor separator gives a global generator}

Suppose the scalar search has produced $x\in C$ and $w\in C\otimes_k C$
with
\[
 (x\otimes1-1\otimes x)w=1-\varepsilon.
\]
Let $f$ be the minimal monic annihilator of $x$, of degree $m>0$.
Put
\[
 H(X,Y)=\frac{f(X)-f(Y)}{X-Y}.
\]
This notation denotes a polynomial constructed by Horner recursion;
it is not division in the tensor algebra. Since $f(x)=0$, the separator
equation and diagonal equalisation give
\[
 H(x\otimes1,1\otimes x)
   =\varepsilon H(x\otimes1,1\otimes x)
   =\varepsilon(1\otimes f'(x)).
\]
Minimality makes $x^{m-1}$ independent of the preceding powers.
The existing Gaussian quotient constructs a $k$-linear functional
$\lambda:C\to k$ which is zero on $1,x,\ldots,x^{m-2}$ and takes
$x^{m-1}$ to $1$. The polynomial $H$ is monic of degree $m-1$ in its
first variable. Contracting that tensor factor therefore gives
\[
 1=f'(x)(\lambda\otimes1)(\varepsilon).
\]
This explicitly constructs the derivative inverse. No regularity is
needed for this step.

Let $Q=k[T]/(f)$. Minimality embeds $Q$ in $C$. It is finite and
reduced, hence regular by the checked theorem. An injective map from a
regular ring reflects units: if $a=a^2b$, the complement $1-ab$
annihilates $a$, and its image must vanish when the image of $a$ is a
unit. Thus $f'(x)^{-1}$ belongs to $Q$.

The coordinates of $H(x\otimes1,1\otimes x)$ in the left tensor basis
belong to $Q$. The formal proof constructs their preimages by Horner
recursion using the multiplication matrix of $x$. The same is then
true of the coordinates $c_j$ of $\varepsilon$, since
$\varepsilon=(1\otimes f'(x)^{-1})H$. Finally
\[
 a=\sum_j\Tr(ab_j)c_j
\]
constructs a preimage in $Q$ of every $a\in C$. Consequently $Q\to C$
is an isomorphism. Degree zero is handled by the equation $1=0$.
This argument replaces selected relative bases and further idempotent
refinements in the large-field branch.

\subsection{All small-field factors, with finite termination}

In the other branch, the actual finite coefficient field $k_0$ and
the lifted monic equations construct a finite reduced subalgebra
$A=k_0[b_1,\ldots,b_n]\subseteq C$. Enumerate all elements of $A$.
For each element $a$ choose a regular witness $a=a^2b$, inside the
finite truncation, and split every current piece by $ab$ and $1-ab$.
On one piece $a$ is invertible; on the other it is zero. Previously
resolved equations survive further refinement. Induction on the
enumeration length supplies termination, without dimension descent.

Every nonzero final corner $eA$ is thus an explicitly enumerated finite
field. Its checked primitive-element construction represents each
$eb_i$ by a polynomial over $k_0$. These representations span $eC$
over $k$ and make it monogenic. Trace nondegeneracy passes to $eC$
and then to its minimal-polynomial quotient, giving separability by
the checked monic trace theorem. Zero pieces are included explicitly.
The isomorphisms
$eC\simeq efC\times e(1-f)C$ assemble the resulting finite product.
Both sides of every split are retained; no prime is chosen.

\section{Checked milestones and remaining local integration}

\begin{enumerate}
\item \textbf{The diagonal and regularity.}
Checked: the presented diagonal, $J=J^2$ directly from
\texttt{Unique}, the determinant construction, general flat base change
and averaging, and $a=a^2b$ over a discrete field.  The exported
\texttt{UnramifiedFieldRegular} theorem also proves reducedness and
algebraic zero-dimensionality.  It assumes only the stated field,
finite-presentation and formal-unramifiedness data, with no assumed
transfer principle or ZMT.
\item \textbf{Finite basis.}
The finite-span part of Proposition~\ref{finite} is checked, using monic
substitutions, induction on algebra generators and integral descent.
This already supplies the finite-fibre input needed before ZMT.
The separate basis construction is also checked, using bounded monic
quotients and quotient Gaussian elimination. It does not require a general
dimension theory or a separate Gr\"obner-basis development.
\item \textbf{One separable monogenic neighbourhood.}
The trace identity is checked with a finite basis as input.
The selected relative-basis, simple-root, small-field and large-field
constructions are checked, together with their bounded-coefficient
dichotomy. The final theorem
\texttt{UnramifiedFieldSelected.Over.At.theorem} takes only a discrete
nonzero field, a finitely presented formally unramified algebra, and a
specified decidable proper prime. It merely supplies a finite basis
and the exact \texttt{SelectedFactorStatement.Selected} certificate.
It assumes neither ZMT nor a field structure theorem.
\item \textbf{The full product structure.}
\texttt{UnramifiedFieldProduct} now gives the product above from the
same field, finite-presentation and formal-unramifiedness hypotheses,
with no prime input. A fresh check gives $18757$ declarations,
$0$ failed, in $9.42$\,s; $847$ declarations in $25$ modules are new.
All $446$ modules pass \texttt{segc all}. Regressions include the base
field $\mathbb F_2$, the constructed standard algebra for $T^2-T$,
the actual large-field separator, $\mathbb F_2^3$ through the small
branch, and the zero algebra. The product theorem is not yet wired
into a replacement for the existing selected-factor local adapter.
\item \textbf{The residual bridge.}
Base-change \texttt{Unique} to $R/\mathfrak m$, transport the chosen
factor through the ZMT localisation, and lift its finitely many
coefficients and generation equations under truncation.  Use the
weaker finite-generator interface described above.  Check the
nonreduced source $B/\mathfrak mB$ and the empty/zero cases explicitly.
In particular, do not assume $B/\mathfrak mB$ injects into the
localised special fibre: each equality pulled back from that fibre
must carry an explicit power of the localisation element $s$.
The base change, the smaller local interface and the generic clearing
lemmas now check. \texttt{QuotientPowerClearing} proves
\[
 \eta(b)=0\text{ in }(B/I)[1/s]
 \quad\Longleftrightarrow\quad
 \exists N\;s^N b\in I.
\]
If $s^N(1-sv)\in I$, it also constructs mere representatives in $B$
for every element of the localisation: replace $b/s^n$ by $bv^n$.
This gives surjectivity, with no claim of injectivity.

\texttt{LocalResidueField} constructs the discrete nonzero field
$R/\mathfrak m$. \texttt{QuotientFibreBaseChange} equips the actual
quotient $E/\mathfrak mE$ with the mapped finite presentation and transfers
formal unramifiedness to it. \texttt{FiniteSpecialFibre} consequently
supplies the usual quotient-module finiteness hypothesis of ZMT.
The assembly of a supplied selected factor's finitely many coefficients
and identities into the exact local interface now checks in
\path{sag/etale-field-selected-lift}. It still takes the selected field
certificate, compatible prime, quasi-inverse relation and local ZMT
data explicitly. Its split-algebra regression includes a nonzero
quotient class which vanishes after localisation.
The subsequent adapter in \path{sag/etale-field-local-complete} now
uses \texttt{UnramifiedFieldSelected} to construct the basis and selected
certificate. It retains explicit finite presentation and formal
unramifiedness of the exact localized fibre, the field structure of the
actual residue quotient, and the compatible prime, quasi-inverse and
local ZMT data.
\end{enumerate}

The existing tensor product of presented algebras
(\texttt{sag/tensor/Tensor.seg}), effective quotients, determinant
trick, flat retractions, principal base change, FieldCFN and HilbertCFN
are reused throughout the checked construction. The existing local
theorem is unchanged. Its final unconditional application still requires
the connections supplying the compatible prime, quasi-inverse relation
and local ZMT data on the actual localized special fibre.

\begin{thebibliography}{9}
\bibitem{decisif} H. Lombardi, \emph{Un th\'eor\`eme d\'ecisif},
2025, \href{https://arxiv.org/abs/2506.07098}{arXiv:2506.07098}.
Local source: \texttt{CANONICITY/decisif.pdf}; especially Theorems 4.3--4.5.
\bibitem{LQ} H. Lombardi and C. Quitt\'e,
\emph{Commutative Algebra: Constructive Methods}, 2015.
Local source: \texttt{COLL/LQ.pdf}; III--9, IV--8.13--8.16,
VI--1.4--1.9, VII--1.4--1.5.  The numbering here is that of this
English PDF.
\bibitem{trace} C. Quitt\'e and H. Lombardi,
\emph{Identit\'es alg\'ebriques permettant de d\'emontrer qu'une
alg\`ebre libre finie nette est traciquement \'etale}, 2025,
\href{https://arxiv.org/abs/2506.03851}{arXiv:2506.03851}, Theorem 3.1.
\bibitem{diagonal} The Stacks Project,
\href{https://stacks.math.columbia.edu/tag/02FL}{Tag 02FL},
the diagonal idempotent for an unramified algebra.
\end{thebibliography}
\end{document}
